%% kapa: exam=1 round=21 year=2010 period=2 form=A subject=appraisal_law subject_ko=부동산관계법규 seq=40 number=80
%% source: Q-Net 감정평가사 기출문제 (cst00309) artlSeq=5024181 file=extracted/1008/5024181/2010년 제21회 감정평가사 1차시험 기출문제/2교시.hwp
%% conversion: kordoc 4.24.51 -> tools/gen1.py (auto); review=proofread-sonnet; flags=none
\begin{question}{80}
부동산등기법령상 등기절차에 관한 설명으로 옳지 않은 것은?
\begin{choices}
  \item 가등기는 신청서에 가등기의무자의 승낙서 또는 가처분명령의 정본을 첨부하여 가등기권리자가 신청할 수 있다.
  \item 등기명의인의 표시의 변경 또는 경정의 등기는 등기명의인 단독으로 신청할 수 있다.
  \item 상속으로 인한 등기는 등기권리자 단독으로 신청할 수 있다.
  \item 국가나 지방자치단체가 소유한 부동산에 관한 등기는 등기권리자의 청구에 의하여 관공서가 지체 없이 촉탁서에 등기원인을 증명하는 서면을 첨부하여 등기소에 촉탁하여야 한다.
  \item 동일한 등기소의 관할 내에 있는 여러 개의 부동산에 관한 등기를 신청하는 경우에는 등기원인이 다를 때에도 동일한 신청서로 등기를 신청할 수 있다.
\end{choices}
\end{question}
\answer{5}
